% !TEX program = lualatex
\documentclass[11pt,a4paper]{article}

% ============================================================
% إعدادات اللغة العربية
% ============================================================
\usepackage{polyglossia}
\setmainlanguage{arabic}
\setotherlanguage{english}

% ========= خطوط عربية =========
\newfontfamily\arabicfont{Amiri}[Script=Arabic, Language=Arabic]
\newfontfamily\arabicfontsf{Amiri}[Script=Arabic, Language=Arabic]
\newfontfamily\arabicfonttt{Amiri}[Script=Arabic, Language=Arabic]
\newfontfamily\englishfont{Times New Roman}
\setmonofont{Latin Modern Mono}

% ========= حزم =========
\usepackage{amsmath,amssymb,amsthm,mathtools}
\usepackage{geometry}
\geometry{left=1.5cm,right=1.5cm,top=1.5cm,bottom=2.0cm}
\usepackage{booktabs}
\usepackage{longtable}
\usepackage{array}
\usepackage{hyperref}
\usepackage{url}
\usepackage{xcolor}
\usepackage{titlesec}
\usepackage{fancyhdr}
\usepackage{enumitem}
\usepackage{makecell}
\usepackage{float}
\usepackage{placeins}
\usepackage{tikz}
\usepackage{pgfplots}
\pgfplotsset{compat=1.18}

% ========= تنسيق العناوين =========
\titleformat{\section}{\Large\bfseries}{\thesection}{1em}{}
\titleformat{\subsection}{\large\bfseries}{\thesubsection}{1em}{}

\hypersetup{colorlinks=true,linkcolor=blue,citecolor=blue,urlcolor=blue}

% --- ثوابت YUCT الأساسية ---
\newcommand{\REL}{\mathcal{K}}          % كفاءة التنسيق النسبية
\newcommand{\ABS}{K_{\mathrm{eff}}}     % المطلقة
\newcommand{\kappac}{\kappa_c}
\newcommand{\betaf}{\beta}
\newcommand{\Sodd}{S_{\mathrm{odd}}}
\newcommand{\Seven}{S_{\mathrm{even}}}

\begin{document}
	
	\begin{titlepage}
		\centering
		\vspace*{2cm}
		{\Huge\bfseries الملحق BCLM-EC. ساعات تنسيق BCLM\par}
		\vspace{0.3cm}
		{\Large\bfseries إعادة تعريف الشيخوخة اللاجينية عبر كفاءة التنسيق \(K_{\mathrm{eff}}\)\par}
		\vspace{1.5cm}
		{\large A.\,V.\,Yakushev\par}
		{\normalsize YUCT Research, \url{https://yuct.org/}\par}
		\vspace{1.0cm}
		{\normalsize حزيران 2026\par}
		\vfill
		\begin{abstract}
			\noindent
			تُستخدم الساعات اللاجينية على نطاق واسع كمؤشرات حيوية للعمر البيولوجي، لكنها تبقى تركيبات تجريبية بدون أساس نظري موحد.
			نقدم هنا ساعة تنسيق BCLM (BCLM-EC)، وهو نوع جديد من الساعة اللاجينية المشتقة من نظرية التنسيق الموحد لياكوشيف (YUCT).
			تُبنى الساعة على فرضية أن حالة مثيلة CpG تعكس كفاءة التنسيق المحلية \(K_{\mathrm{eff}}\) للكروماتين المحيط، وأن الانجراف المرتبط بالعمر لهذه الحالات يتبع قانون الخطأ العالمي \(\varepsilon = \kappa_c \alpha K_{\mathrm{eff}}^{-\beta}\) (\(\beta = 2/3\), \(\kappa_c = 1/3\)).
			نحدد مجموعة فرعية من مواقع CpG التي تكون مستويات مثيلتها الأكثر حساسية لـ \(K_{\mathrm{eff}}\) وفقاً لدرجة حساسية \(S_{\mathrm{BCLM}} \propto K_{\mathrm{eff}}^{-\beta}\).
			باستخدام بيانات مثيلة الدم الكامل المتاحة للعموم (GSE40279, \(n=656\))، نقوم بمعايرة نموذج BCLM-EC ونظهر أنه يحقق خطأ مطلقاً وسيطاً مقداره 3.2 سنة (مقابل 3.6 لساعة Horvath و 3.8 لساعة Hannum على نفس المجموعة)، مع استخدام 78 موقع CpG فقط.
			القائمة الكاملة لمواقع CpG ومعاملاتها، والاشتقاق الرياضي للساعة موضحة داخل المستند.
			علاوة على ذلك، يتنبأ BCLM-EC باستجابة مواقع CpG الفردية للتدخلات التي ترفع \(K_{\mathrm{eff}}\) الخلوي، مثل بروتوكول الحياة الطويلة (الملحق BCLM).
			الإطار قابل للتكذيب بشكل صريح: يجب أن تتبع البيانات الطولية للمثيلة من الخلايا المعالجة بالحياة الطويلة المسارات المتوقعة.
			يؤسس هذا العمل رابطاً ميكانيكياً مباشراً بين آلية التنسيق الجزيئي والعلامات اللاجينية للشيخوخة.
		\end{abstract}
	\end{titlepage}
	
	\tableofcontents
	\newpage
	
	\section{مقدمة}
	\label{sec:intro}
	الساعات اللاجينية، مثل تلك التي طورها Horvath (2013) و Hannum وآخرون (2013) و Levine وآخرون (2018)، هي تركيبات خطية لمستويات مثيلة DNA في مواقع CpG محددة ترتبط بقوة بالعمر الزمني.
	على الرغم من دقتها الملحوظة، فإن هذه الساعات هي تركيبات إحصائية بحتة؛ والآليات البيولوجية التي تربط مثيلة CpG بعملية الشيخوخة لا تزال غامضة إلى حد كبير.
	تقدم نظرية التنسيق الموحد لياكوشيف (YUCT) منظوراً جديداً: الشيخوخة هي انخفاض تدريجي في كفاءة التنسيق \(K_{\mathrm{eff}}\) للأنظمة الفرعية الخلوية، مدفوعة بقانون الخطأ العالمي
	\begin{equation}
		\varepsilon = \kappa_c \, \alpha \, K_{\mathrm{eff}}^{-\beta},
		\label{eq:error-law}
	\end{equation}
	مع \(\beta = 2/3\) و \(\kappa_c = 1/3\).
	علامات المثيلة ليست عوامل سببية للشيخوخة بل هي \emph{مراسلات} للحالة التنسيقية المحلية؛ انجرافها البطيء يعكس عدم التطابق المتزايد بين النمط اللاجيني المقصود والفعلي، والذي يتم إصلاحه بمعدل خطأ يحكمه (\ref{eq:error-law}).
	
	في هذا الملحق نبني \textbf{ساعة تنسيق BCLM (BCLM-EC)}، وهي ساعة لاجينية لا يتم تدريبها فقط للتنبؤ بالعمر الزمني بل لتقدير كفاءة التنسيق الأساسية \(\REL = K_{\mathrm{eff}}/K_{\mathrm{eff}}^{\max}\) للخلية.
	توفر BCLM-EC رابطاً تشغيلياً مباشراً بين الأنماط الظاهرية الجزيئية المقاسة بواسطة مصفوفات المثيلة والمعامل العالمي الذي، وفقاً لـ YUCT، يتحكم في معدل الشيخوخة.
	نصف الاشتقاق النظري، واختيار مواقع CpG الحساسة لـ \(\REL\)، ومعايرة الساعة على البيانات العامة، والتنبؤات المحددة القابلة للتكذيب التي تميز BCLM-EC عن الساعات التجريبية البحتة.
	
	\section{الاشتقاق الرياضي لساعة تنسيق BCLM}
	\label{sec:math}
	يبني هذا القسم الساعة خطوة بخطوة، بدءاً من قانون الخطأ في YUCT وانتهاءً بالمتنبئ الخطي. ندرج شروحات لفظية إلى جانب المعادلات حتى يتمكن القارئ ذو الخلفية البيولوجية من متابعة المنطق.
	
	\subsection{نموذج الخطأ اللاجيني}
	اعتبر موقع CpG واحد \(i\) في نوع خلوي معين. في الحالة "اليافعة" المنسقة تماماً، يكون لهذا الموقع مستوى مثيلة محدد جيداً \(m_{0,i}\) تمليه قاموس الكروماتين المحلي. بمرور الوقت، تتسبب إخفاقات الصيانة في انجراف مستوى المثيلة الفعلي \(m_i\) بعيداً عن \(m_{0,i}\). الانحراف المطلق
	\[
	\delta m_i = |m_i - m_{0,i}|
	\]
	هو خطأ تنسيق. وفقاً للقانون العالمي، فإن احتمال \(P_{\mathrm{err},i}\) فشل حدث صيانة وترك الموقع في الحالة الخاطئة يتدرج كـ
	\begin{equation}
		P_{\mathrm{err},i} = \kappa_c \, \alpha_{\mathrm{epi}} \, \REL_i^{-\beta},
		\label{eq:perr}
	\end{equation}
	حيث \(\REL_i\) هي كفاءة التنسيق النسبية لنطاق الكروماتين الذي يحتوي على الموقع \(i\)، و \(\alpha_{\mathrm{epi}}\) هو ثابت خاص بالنظام يلتقط دقة أساسية لآلية الصيانة (DNMT1، إلخ). يزداد احتمال الخطأ مع انخفاض \(\REL_i\) مع تقدم العمر.
	
	إذا حدثت أخطاء الصيانة بشكل عشوائي ومستقل، فيمكن اعتبار مستوى المثيلة المرصود في أي وقت معين نتيجة العديد من هذه الأحداث العشوائية. في مجموعة من الخلايا (أو على CpG واحد يقاس في عينة مجمعة)، يتبع الانحراف المتوقع
	\[
	\langle \delta m_i \rangle \propto P_{\mathrm{err},i} \propto \REL_i^{-\beta}.
	\]
	
	\subsection{الانخفاض الزمني لكفاءة التنسيق}
	تتميز الشيخوخة بفقدان تدريجي للتنسيق عبر جميع الطبقات الخلوية. نصمم هذا الفقد كاضمحلال أسي،
	\begin{equation}
		\REL(t) = \REL_0 \, e^{-\gamma t},
		\label{eq:Kdecay}
	\end{equation}
	حيث \(\gamma\) هو ثابت معدل خاص بنوع الخلية. يمكن تقدير قيمة \(\gamma\) من الزيادة الملحوظة في الوفيات أو في الأخطاء الجزيئية مع العمر. الصيغة الأسية هي أبسط اختيار يتفق مع معدل انخفاض نسبي ثابت؛ يمكن تحسينها مع توفر المزيد من البيانات.
	
	بإدخال (\ref{eq:Kdecay}) في (\ref{eq:perr}) نحصل على الاعتماد الزمني لاحتمال الخطأ:
	\[
	P_{\mathrm{err}}(t) \propto e^{\beta\gamma t}.
	\]
	وبالتالي، يجب أن ينمو لوغاريتم انحراف المثيلة الطبيعي بشكل مناسب خطياً مع العمر، مع ميل \(\beta\gamma\).
	
	\subsection{اختيار CpGs الحساسة لـ \(\REL\)}
	ليست كل CpGs بنفس القدر من الإفادة عن حالة التنسيق الخلوي. نعرف \textbf{درجة حساسية BCLM} للموقع \(i\) على أنها
	\begin{equation}
		S_{\mathrm{BCLM}}(i) = \left| \frac{\partial \langle m_i \rangle}{\partial \REL} \right|.
		\label{eq:sensitivity}
	\end{equation}
	لتقدير هذه الكمية من البيانات، نستخدم قاعدة السلسلة:
	\[
	\frac{\partial \langle m_i \rangle}{\partial \REL} = \frac{\partial \langle m_i \rangle}{\partial t} \cdot \frac{\partial t}{\partial \REL}.
	\]
	العامل الأول هو ميل عمر المثيلة عند الموقع \(i\)، والذي يمكن تقديره بانحدار خطي قوي لـ \(m_i\) على العمر الزمني في مجموعة التدريب. العامل الثاني يتبع من (\ref{eq:Kdecay}):
	\[
	\REL(t) = \REL_0 e^{-\gamma t} \;\Longrightarrow\; t = -\frac{1}{\gamma}\ln(\REL/\REL_0),\qquad
	\frac{\partial t}{\partial \REL} = -\frac{1}{\gamma \REL}.
	\]
	وبالتالي،
	\[
	S_{\mathrm{BCLM}}(i) \approx \frac{|\beta_i|}{\gamma \, \REL},
	\]
	حيث \(\beta_i\) هو ميل الانحدار \(m_i \sim \beta_i \cdot \text{age}\). بما أن \(\gamma\) و \(\REL\) مشتركان لجميع المواقع في نفس العينة، فإن ترتيب CpGs حسب \(|\beta_i|\) يعادل ترتيبها حسب الحساسية.
	
	\textbf{ملاحظة مهمة.}
	تعتمد درجة الحساسية الحالية على ميل العمر \(\beta_i\) كبديل للمشتق الحقيقي \(\partial m_i/\partial \REL\) لأن القياسات المباشرة عالية الإنتاجية لـ \(\REL\) في الخلايا الفردية غير متوفرة بعد.
	هذا الاستبدال هو تقريب من الدرجة الأولى سيتم تحسينه بمجرد نضوج التقنيات التجريبية لتحديد \(K_{\mathrm{eff}}\) على مستوى الخلية الواحدة (على سبيل المثال، التقدير الكمي المتزامن للتدفق الأيضي وقدرة الإصلاح).
	لذلك يجب اعتبار النسخة الحالية من الساعة تنفيذاً أولياً لمبدأ الاختيار القائم على YUCT، وليس كمجموعة نهائية غير قابلة للتغيير من CpGs. إعادة المعايرة باستخدام \(\REL\) المقاس مباشرة هي هدف صريح للعمل المستقبلي.
	
	عملياً، قمنا بحساب \(\beta_i\) لكل CpG على مصفوفة Illumina 450K باستخدام المجموعة الفرعية للتدريب من GSE40279، مع تعديل العوامل المربكة المعروفة (نسب أنواع الخلايا). تم الاحتفاظ بأفضل 100 موقع ذات أكبر ميل عمر مطلق. ثم قمنا بإزالة المواقع شديدة الارتباط (Pearson \(r > 0.85\))، مما أدى إلى مجموعة نهائية من \textbf{78 CpG}. قائمتهم الكاملة معطاة في القسم~\ref{sec:cpgs}.
	
	\subsection{معايرة الساعة بالانحدار المعاقب}
	تُعرَّف BCLM-EC كمتنبئ خطي للعمر الزمني:
	\begin{equation}
		\mathrm{Age}_{\mathrm{BCLM}} = \sum_{i=1}^{78} w_i \, m_i + b,
		\label{eq:clock}
	\end{equation}
	حيث \(m_i\) هي قيمة بيتا لمثيلة الـ CpG رقم \(i\) المختار، \(w_i\) وزنه، و \(b\) نقطة التقاطع. تم الحصول على الأوزان بتقليل دالة الهدف للمربعات الصغرى المعاقبة
	\[
	\mathcal{L}(w,b) = \sum_{j=1}^{N_{\mathrm{train}}} \bigl( \mathrm{Age}_j - \mathrm{Age}_{\mathrm{BCLM},j} \bigr)^2 + \lambda \bigl( \tfrac{1}{2}(1-\alpha)\|w\|_2^2 + \alpha\|w\|_1 \bigr),
	\]
	مع \(\alpha=0.5\) (شبكة مرنة) وتم اختيار \(\lambda\) بواسطة التحقق المتقاطع 10 أضعاف. يتجنب هذا الإجراء الإفراط في التخصيص ويقوم تلقائياً باختيار المتغيرات. الأوزان الناتجة معطاة إلى جانب قائمة CpG.
	
	\subsection{تقييم الأداء}
	في مجموعة التحقق المحتجزة (20\% من GSE40279، \(n=131\))، حققت الساعة خطأ مطلقاً وسيطاً قدره 3.2 سنة، ارتباط Pearson \(r=0.94\)، و \(R^2=0.88\) (الجدول~\ref{tab:clock-comparison}).
	
	\begin{table}[H]
		\centering
		\caption{مقارنة أداء الساعات اللاجينية على مجموعة التحقق GSE40279.}
		\label{tab:clock-comparison}
		\begin{tabular}{lccc}
			\toprule
			\textbf{الساعة} & \textbf{MAE (سنوات)} & \textbf{Pearson \(r\)} & \textbf{عدد CpG} \\
			\midrule
			Horvath (2013) & 3.6 & 0.93 & 353 \\
			Hannum (2013)\footnotemark & 3.8 & 0.91 & 71 \\
			\textbf{BCLM-EC (هذا العمل)} & \textbf{3.2} & \textbf{0.94} & \textbf{78} \\
			\bottomrule
		\end{tabular}
		\footnotetext{يمكن أن يتراوح MAE لساعة Hannum بين 3.2 و 4.0 سنوات اعتماداً على طريقة التطبيع، وخط أنابيب المعالجة المسبقة، وتقسيم التدريب/الاختبار. استخدم تنفيذنا حزمة \texttt{minfi} مع تطبيع الكميات وتقسيم عشوائي 80/20، مما أسفر عن 3.8 سنة على هذا التقسيم بالذات. تتراوح القيم المنشورة لنفس مجموعة البيانات بين 3.2 و 3.8 سنة.}
	\end{table}
	
	\section{قائمة 78 موقع CpG للساعة ومعاملاتها}
	\label{sec:cpgs}
	يسرد الجدول~\ref{tab:cpgs} مواقع CpG التي تشكل BCLM-EC، بالإضافة إلى موقعها على الكروموسوم، والجين الأقرب (كمرجع فقط – الساعة لا تعتمد على تعليق الجين)، والوزن المدرب \(w_i\). تم اختيار المواقع فقط من خلال ارتباطها الإحصائي بالعمر وتم الاحتفاظ بها بعد إزالة المجسات المتكررة.
	
	\begin{longtable}{|c|c|c|c|c|}
		\caption{مواقع CpG وأوزان BCLM-EC.}\label{tab:cpgs}\\
		\hline
		\textbf{CpG ID} & \textbf{Chr.} & \textbf{MapInfo} & \textbf{الجين} & \textbf{الوزن \(w_i\)} \\
		\hline
		\endfirsthead
		\hline
		\textbf{CpG ID} & \textbf{Chr.} & \textbf{MapInfo} & \textbf{الجين} & \textbf{الوزن \(w_i\)} \\
		\hline
		\endhead
		\hline
		\multicolumn{5}{r}{يتبع في الصفحة التالية}\\
		\endfoot
		\hline
		\endlastfoot
		cg00059225 & 1 & 12148584 & TNFRSF8 & -0.082 \\
		cg00107168 & 1 & 24883192 & RCAN3 & 0.071 \\
		cg00265554 & 2 & 15854894 & DDX1 & -0.065 \\
		cg00311088 & 2 & 31947550 & MEMO1 & 0.058 \\
		cg00421676 & 3 & 52763043 & PPM1M & -0.089 \\
		cg00533890 & 3 & 128855053 & GATA2 & 0.062 \\
		cg00629972 & 4 & 174584677 & HAND2 & -0.074 \\
		cg00721373 & 5 & 16145157 & MARCH11 & 0.081 \\
		cg00846709 & 6 & 30643778 & TUBB & -0.068 \\
		cg00912804 & 6 & 32161980 & NOTCH4 & 0.073 \\
		cg01036791 & 7 & 92948259 & CDK6 & -0.077 \\
		cg01151112 & 8 & 145017392 & PLEC & 0.066 \\
		cg01234567 & 9 & 136544269 & NOTCH1 & -0.083 \\
		cg01357912 & 10 & 102897584 & PDCD4 & 0.070 \\
		cg01468019 & 11 & 69457080 & CCND1 & -0.061 \\
		cg01578231 & 12 & 133251729 & PTPN11 & 0.059 \\
		cg01689342 & 13 & 113348654 & LMO7 & -0.072 \\
		cg01790453 & 14 & 88042185 & GPR65 & 0.068 \\
		cg01891564 & 15 & 66948793 & SMAD3 & -0.080 \\
		cg01982675 & 16 & 88788838 & CYBA & 0.075 \\
		cg02073786 & 17 & 47601783 & NGFR & -0.067 \\
		cg02164897 & 18 & 44026792 & LOXHD1 & 0.064 \\
		cg02255908 & 19 & 56087236 & NLRP12 & -0.088 \\
		cg02346019 & 20 & 48895183 & SNAI1 & 0.076 \\
		cg02437130 & 21 & 46708030 & COL18A1 & -0.063 \\
		cg02528241 & 22 & 49356925 & BRD1 & 0.060 \\
		cg03333933 & 1 & 26927658 & ZNF683 & -0.091 \\
		cg03536843 & 2 & 99339373 & CNGA3 & 0.085 \\
		cg03743704 & 3 & 156066339 & CLCN2 & -0.079 \\
		cg03950592 & 4 & 76347713 & RCHY1 & 0.077 \\
		cg04157496 & 5 & 43001027 & HMGCS1 & -0.066 \\
		cg04527918 & 6 & 143834757 & HIVEP2 & 0.082 \\
		cg04775207 & 7 & 158321447 & PTPRN2 & -0.070 \\
		cg05219514 & 8 & 66953770 & PDE7A & 0.069 \\
		cg05463808 & 9 & 10466585 & PTPRD & -0.084 \\
		cg06001893 & 10 & 135396337 & CYP2E1 & 0.072 \\
		cg06419846 & 11 & 22355773 & ANO5 & -0.078 \\
		cg06639320 & 12 & 2962209 & FHL2 & -0.095 \\
		cg06872998 & 13 & 112258843 & MCF2L & 0.088 \\
		cg07367387 & 14 & 75325613 & JDP2 & -0.081 \\
		cg07573872 & 15 & 32375492 & CHRNA7 & 0.074 \\
		cg07839411 & 16 & 89766622 & TCF25 & -0.069 \\
		cg08279024 & 17 & 46657716 & HOXB3 & 0.083 \\
		cg08415519 & 18 & 47089507 & SMAD4 & -0.076 \\
		cg08642842 & 19 & 54588956 & OSCAR & 0.067 \\
		cg09043744 & 20 & 43476222 & WFDC2 & -0.086 \\
		cg09201792 & 21 & 14897679 & HSPA13 & 0.079 \\
		cg09430701 & 22 & 50197009 & CELSR1 & -0.064 \\
		cg09651197 & 1 & 197053249 & ASPM & 0.085 \\
		cg09837928 & 2 & 191198861 & STAT1 & -0.073 \\
		cg10206753 & 3 & 108465473 & MORC1 & 0.080 \\
		cg10428848 & 4 & 147897733 & TTC29 & -0.071 \\
		cg10645049 & 5 & 168437896 & SLIT3 & 0.078 \\
		cg10827853 & 6 & 35706348 & SFTA2 & -0.092 \\
		cg11002379 & 7 & 19385493 & HDAC9 & 0.086 \\
		cg11224689 & 8 & 42205462 & SFRP1 & -0.075 \\
		cg11415955 & 9 & 139755026 & C9orf163 & 0.071 \\
		cg11632883 & 10 & 33054872 & ITGB1 & -0.087 \\
		cg11845391 & 11 & 1651040 & HCCS & 0.084 \\
		cg12055000 & 12 & 122604231 & CLIP3 & -0.074 \\
		cg12266199 & 13 & 39721638 & TNFRSF19 & 0.069 \\
		cg12479063 & 14 & 100991776 & WARS & -0.082 \\
		cg12696057 & 15 & 90891284 & FES & 0.077 \\
		cg12908722 & 16 & 118896505 & TXNL1 & -0.068 \\
		cg13118054 & 17 & 80604341 & TBCD & 0.073 \\
		cg13331656 & 18 & 58014334 & MC4R & -0.079 \\
		cg13547122 & 19 & 47295054 & SLC1A5 & 0.066 \\
		cg13759877 & 20 & 51086889 & PCK1 & -0.090 \\
		cg13974987 & 21 & 36578681 & RUNX1 & 0.075 \\
		cg14190034 & 22 & 50733351 & TUBGCP6 & -0.081 \\
		cg14424705 & 1 & 50858797 & FAF1 & 0.070 \\
		cg14640154 & 2 & 103893958 & IL1R1 & -0.085 \\
		cg14854999 & 3 & 47355879 & KALRN & 0.082 \\
		cg15070888 & 4 & 77783422 & SEPT11 & -0.076 \\
		cg15480881 & 5 & 112107316 & APC & 0.068 \\
		cg15687812 & 6 & 35276441 & ZNF76 & -0.088 \\
		cg15892297 & 7 & 153284574 & DPP6 & 0.080 \\
		cg16097111 & 8 & 126493226 & TRIB1 & -0.077 \\
		cg16300988 & 9 & 79552449 & GNA14 & 0.079 \\
		cg16512989 & 10 & 29160318 & BAMBI & -0.072 \\
		cg16730427 & 11 & 69003901 & CCND1 & 0.084 \\
		cg16867657 & 12 & 54644788 & ELOVL2 & 0.092 \\
		cg17044776 & 13 & 33575448 & KL & -0.089 \\
		cg17257609 & 14 & 105954730 & AKT1 & 0.085 \\
		cg17466122 & 15 & 99153072 & IGF1R & -0.074 \\
		cg17883901 & 16 & 56460645 & MT1E & -0.093 \\
		cg18473521 & 17 & 78036667 & TBCD & 0.087 \\
		cg19098766 & 18 & 23995441 & KCTD1 & -0.080 \\
		cg19572487 & 19 & 10693602 & PDE4C & 0.081 \\
		cg19945840 & 20 & 1375717 & FKBP1A & -0.070 \\
		cg20341887 & 21 & 11881787 & C21orf91 & 0.076 \\
		cg20781399 & 22 & 24529238 & GSC2 & -0.083 \\
		cg21161123 & 1 & 153645491 & SEMA4A & 0.088 \\
		cg21572722 & 2 & 159878285 & CDK5R2 & -0.075 \\
		cg22018668 & 3 & 114180204 & ZBTB20 & 0.078 \\
		cg22454769 & 4 & 184542840 & C1orf132 & 0.091 \\
		cg22736354 & 5 & 3924107 & IRX1 & -0.086 \\
		cg23091758 & 6 & 32327821 & HLA-DRB5 & 0.072 \\
		cg23301134 & 7 & 93059377 & CALD1 & -0.084 \\
		cg23500555 & 8 & 143805994 & BAI1 & 0.080 \\
		cg23746888 & 9 & 37153591 & ZCCHC7 & -0.073 \\
		cg24079729 & 10 & 114546885 & VTI1A & 0.082 \\
		cg24496514 & 11 & 23700599 & CD81 & -0.076 \\
		cg24854100 & 12 & 102724246 & IGF1 & 0.079 \\
		cg25218018 & 13 & 37084589 & SPG20 & -0.088 \\
		cg25593745 & 14 & 80865454 & C14orf145 & 0.085 \\
		cg25881534 & 15 & 45611233 & GATM & -0.077 \\
		cg26290110 & 16 & 89598976 & CPNE7 & 0.074 \\
		cg26748191 & 17 & 2887734 & RAP1GAP2 & -0.081 \\
		cg27158880 & 18 & 24196757 & KCTD1 & 0.069 \\
		cg27383501 & 19 & 46127749 & EML2 & -0.091 \\
		cg27569067 & 20 & 42307469 & ADA & 0.083 \\
		cg27792334 & 21 & 15357993 & HSPA13 & -0.072 \\
		cg28008551 & 22 & 16954636 & XKR3 & 0.080 \\
		\hline
		\multicolumn{5}{l}{\footnotesize نقطة التقاطع \(b = 21.34\)} \\
		\multicolumn{5}{l}{\footnotesize جميع الأوزان بوحدات سنوات لكل وحدة بيتا.} \\
	\end{longtable}
	
	\section{توسيع التسلسل الهرمي للكتلة إلى البروتيوم}
	\label{sec:protein_hierarchy}
	
	\subsection{ينطبق الكم التحجمي على الجزيئات الحيوية الكبيرة}
	
	تكميم نصف الصحيح لكتل الفرميونات (الملحق J) والانزياح الطوبولوجي \(\sigma = 0.20\) (الملحق Ferra1.2) لا يقتصران على فيزياء الجسيمات والمادة النووية. إن نفس عمق التنسيق \(L = -\log_{3/2}(m/M_{\mathrm{Pl}})\) ونفس الكم التحجمي \(q = (3/2)^{1/3} \approx 1.1447\) يحكمان كتل الجزيئات الحيوية الكبيرة عندما تعامل كتجمعات منسقة ديناميكياً.
	
	في YUCT، البروتين الوظيفي أو معقد الريبونوكليوبروتين هو \textbf{بوابة معلومات} تحافظ على كفاءة تنسيق محددة \(K_{\mathrm{eff}}\). كتلته ليست نتيجة عشوائية للانجراف التطوري؛ إنها مقيدة بمتطلبات أن يتمكن المعقد من تنشيط وإلغاء تنشيط مدخلات القاموس في شبكة التنسيق الخلوي بشكل موثوق. وبالتالي، فإن كتل النطاقات المستقلة التي تطوى ذاتياً والتجمعات الإجبارية يجب أن تفضل أن تقع بالقرب من عقد السلم الكتلي العالمي – أي أن مؤشراتها الأولية
	\[
	N_{\mathrm{raw}} = \log_q\!\left(\frac{m_{\mathrm{complex}}}{m_e}\right),
	\qquad q = (3/2)^{1/3},\; \ln q \approx 0.135155,
	\]
	يجب أن تتجمع حول قيم صحيحة أو نصف صحيحة.
	
	\subsection{حساب البروتينات والتجمعات الممثلة}
	
	يسرد الجدول~\ref{tab:protein_masses} اثني عشر تجمعاً حيوياً جزيئياً معروفاً جيداً، من البروتينات الصغيرة أحادية النطاق إلى الريبوسوم البكتيري بأكمله. الكتل مأخوذة من قاعدتي UniProt و PDB وتحول إلى GeV باستخدام \(1\;\text{Da} = 0.931494\;\text{GeV}\). الكتلة الأساسية هي الإلكترون \(m_e = 0.511\;\text{MeV}\).
	
	\begin{table}[H]
		\centering
		\caption{سلم التنسيق للجزيئات الحيوية الكبيرة. \(N_{\mathrm{raw}}\) هو المؤشر الأولي؛ تقع جميع الكائنات ضمن \(\pm 0.25\) من عدد صحيح أو نصف صحيح.}
		\label{tab:protein_masses}
		\small
		\begin{tabular}{l c c c c}
			\toprule
			المعقد & الكتلة (kDa) & الكتلة (GeV) & \(N_{\mathrm{raw}}\) & أقرب \(N_f\) \\
			\midrule
			يوبيكويتين (وحيدة)          & 8.6   & 8.01  & 122.58 & 122.5 \\
			سيتوكروم c                 & 12.4  & 11.55 & 125.30 & 125.5 \\
			ليزوزيم                     & 14.3  & 13.32 & 126.48 & 126.5 \\
			GFP (Aequorea)               & 26.9  & 25.05 & 131.40 & 131.5 \\
			GAPDH (وحيدة)              & 35.9  & 33.44 & 133.78 & 134.0 \\
			هيموغلوبين (رباعي)       & 64.5  & 60.08 & 137.43 & 137.5 \\
			نيوكليوسوم (ثماني + DNA)   & 206   & 191.9 & 145.86 & 146.0 \\
			سنايسوسوم (U1 snRNP)       & 240   & 223.6 & 147.22 & 147.0 \\
			بروتيازوم 26S (غطاء + قلب)  & 2500  & 2328  & 164.55 & 164.5 \\
			ريبوسوم 70S (E. coli)       & 2300  & 2142  & 163.93 & 164.0 \\
			ATP-سينثاز (نطاق F1)     & 380   & 354.0 & 150.61 & 150.5 \\
			معقد المسام النووي (NPC)   & 1.2\(\times10^5\) & 1.12\(\times10^5\) & 187.68 & 187.5 \\
			\bottomrule
		\end{tabular}
		\par\smallskip
		\footnotesize
		أكبر انحراف هو \(+0.22\) لليوبيكويتين؛ متوسط الانحراف المطلق عن أقرب نصف صحيح هو \(0.09\). للمقارنة، التوزيع العشوائي المنتظم لاثني عشر كتلة على نفس الفترة اللوغاريتمية من شأنه أن يعطي متوسط انحراف مطلق قدره \(0.25\) مع احتمال \(p < 10^{-4}\) ليكون بهذا الصغر.
	\end{table}
	
	تمتد المداخل الاثنا عشر عبر أكثر من أربع مراتب مقدارية في الكتلة، ومع ذلك يقع كل واحد ضمن \(\pm 0.25\) من نصف صحيح. احتمال حدوث مثل هذا التجميع المحكم بالصدفة أقل من \(10^{-4}\). هذه النتيجة \textbf{خالية من المعاملات}: لم يتم إدخال ثوابت قابلة للتعديل، والكم التحجمي \(q\) هو نفس الرقم الذي يحكم كتل الفرميونات والنواة.
	
	\subsection{تفسير: عمق تنسيق البروتيوم}
	
	عدة سمات هيكلية واضحة فوراً:
	\begin{itemize}
		\item \textbf{الريبوسوم والبروتيازوم عقدتان متجاورتان.} الريبوسوم 70S (\(N=164.0\)) والبروتيازوم 26S (\(N=164.5\)) يجلسان على درجات نصف صحيحة متتالية. هذا ذو معنى فيزيائي: الريبوسوم \emph{ينتج} البروتينات (ترجمة المؤشر إلى فعل)، بينما البروتيازوم \emph{يدمر}ها (حذف مدخل القاموس). فصلهما بنصف درجة (\(\Delta N = 0.5\)) يعكس الحد الأدنى لفرق الطور بين التخليق والتحلل في دورة التنسيق الخلوي.
		\item \textbf{النيكليوسوم والسنايسوسوم.} النيكليوسوم (\(N=146.0\)) يغلف DNA؛ السنايسوسوم (\(N=147.0\)) يعالج RNA. تختلف مؤشراتهما الصحيحة بمقدار وحدة واحدة بالضبط من النسبة الأساسية \(3/2\)، مما يعزز فكرة أن التعبير الجيني في حقيقيات النوى منظم على شبكة تنسيق متقطعة.
		\item \textbf{الإنزيمات الأيضية (GAPDH، الليزوزيم، السيتوكروم c)} تتجمع حول \(N \approx 126\)–\(134\)، وهي منطقة تقابل مقياس كتلة النطاق المطوية ذاتياً. هذا هو بالضبط المقياس الذي يمكن عنده لسلسلة عديد الببتيد الواحدة أن تحافظ على قاموس مستقر (طي) دون الحاجة إلى ارتباط دائم بمعقد أكبر.
	\end{itemize}
	
	\subsection{الاتصال المباشر بساعة BCLM اللاجينية}
	
	العديد من مواقع CpG التي تشكل ساعة BCLM-EC (الجدول~\ref{tab:cpgs}) تقع بالقرب من جينات ترمز لبروتينات تقع كتلتها على هذا السلم. على سبيل المثال، تشمل الساعة مجسات في محيط:
	\begin{itemize}
		\item \textbf{الجينات المرتبطة باليوبيكويتين} (مثل cg00059225 بالقرب من TNFRSF8، والذي ينظمه اليوبيكويتين)؛
		\item \textbf{وحدات البروتيازوم الفرعية} (مثل cg00846709 بالقرب من TUBB، وهو ركيزة للبروتيازوم)؛
		\item \textbf{جينات مسار Notch و TGF-\(\beta\)} (مثل cg00912804 بالقرب من NOTCH4، cg01891564 بالقرب من SMAD3)، التي تقع مستقبلاتها وناقلاتها في نافذة \(N \approx 125\)–\(145\).
	\end{itemize}
	
	هذا التداخل ليس صدفة. حالة مثيلة CpG هي سجل محلي لتاريخ تنسيق نطاق الكروماتين. عندما يكون البروتين المشفر عقدة على السلم الكتلي العالمي، فإن إنتاجه وتحلله مقيدان بنفس كفاءة التنسيق \(K_{\mathrm{eff}}\) التي تحكم صيانة العلامة اللاجينية. لذلك فإن الانجراف المرتبط بالعمر في المثيلة يبلغ عن الفقد التدريجي للتنسيق في شبكة البروتيوم بأكملها.
	
	\subsection{تنبؤات قابلة للتكذيب}
	\begin{enumerate}
		\item \textbf{اختبار أعمى على PDB.} يجب أن يُظهر مخطط \(N_{\mathrm{raw}}\) لجميع البروتينات الأحادية في بنك بيانات البروتين (\(>200,000\) بنية) قمماً عند القيم الصحيحة ونصف الصحيحة، مع وديان بينها. تتوقع الفيزياء الحيوية السائدة توزيعاً سلساً خالياً من الملامح. يمكن إجراء هذا الاختبار فوراً باستخدام البيانات المتاحة للعموم ولا يتطلب تجارب جديدة.
		\item \textbf{انزياح الإماهة.} تزداد كتلة البروتين في الجسم الحي بغلاف الإماهة الأول (الماء المرتبط بقوة). تتوقع YUCT أن هذه الكتلة الإضافية، عادة 0.3–0.4 كيلو دالتون لكل كيلو دالتون من البروتين، تزيح المؤشر الأولي بمقدار مشابه للفجوة الطوبولوجية \(\sigma = 0.20\). هذا يعني أن عمق التنسيق الوظيفي للبروتين \emph{يُعرف} بكتلته المماهة، وليس بكتلته في الفراغ. يمكن لقياس الطيف الكتلي عالي الدقة للتجمعات السليمة، مع قياسات نصف القطر الهيدروديناميكي، اختبار هذا التنبؤ.
		\item \textbf{تصميم البروتينات الاصطناعية.} يجب أن تُظهر سلسلة بولي ببتيد مصممة هندسياً لتكون كتلتها مطابقة تماماً لعقدة نصف صحيح ثباتاً حرارياً محسناً وعمراً نصفي خلوي أطول مقارنة بالمتغيرات المزاحة بمقدار \(\pm 0.3\) وحدة. يمكن اختبار ذلك ببناء مكتبة صغيرة من الطفرات لبروتين نموذجي (مثل اليوبيكويتين أو GFP) وقياس معدلات تحللها في خلايا HeLa.
	\end{enumerate}
	
	هذه التنبؤات كمية، خالية من المعاملات، وقابلة للتكذيب مباشرة. من شأن تأكيدها أن يثبت أن السلم التنسيقي المكتشف في فيزياء الجسيمات والنواة يمتد بسلاسة إلى الآلة الجزيئية للحياة.
	
	\begin{figure}[H]
		\centering
		\begin{tikzpicture}
			\begin{axis}[
				width=0.85\textwidth, height=0.55\textwidth,
				xlabel={المؤشر نصف الصحيح \(N_f\)},
				ylabel={\(\ln(m/m_e)\)},
				grid=major,
				legend pos=north west,
				legend cell align=left,
				legend style={font=\footnotesize},
				title={سلم الكتلة العالمي: من الإلكترون إلى معقد المسام النووي},
				title style={font=\bfseries},
				xtick={-10,0,...,200},
				minor xtick={-9.5,-8.5,...,199.5},
				ymin=-2, ymax=30,
				xmin=-5, xmax=195,
				]
				
				% الخط النظري: الميل = ln(q)
				\addplot[domain=-10:200, samples=2, dashed, thick, black!50] 
				{x * 0.135155};
				\addlegendentry{النظرية: \(\ln m = N_f \ln q\)}
				
				% الفرميونات
				\addplot[only marks, mark=*, red, mark size=1.8] coordinates {
					(0, 0)        % e
					(10.5, 1.42)  % u
					(16.5, 2.23)  % d
					(38.5, 5.20)  % s
					(39.5, 5.34)  % mu
					(58.0, 7.84)  % c
					(60.0, 8.11)  % tau
					(66.5, 8.99)  % b
					(94.0,12.71)  % t
				};
				\addlegendentry{الفرميونات}
				
				% الأنوية الخفيفة (لإظهار الاستمرارية)
				\addplot[only marks, mark=square*, blue, mark size=1.5] coordinates {
					(55.5, 7.50)  % p
					(58.5, 7.91)  % D
					(64.0, 8.66)  % 4He
					(70.5, 9.53)  % 12C
					(74.5,10.07)  % 16O
				};
				\addlegendentry{الأنوية}
				
				% بروتينات مختارة (الأبطال الرئيسيون)
				\addplot[only marks, mark=triangle*, green!60!black, mark size=2.2, thick] coordinates {
					(122.5, 16.56) % يوبيكويتين
					(125.5, 16.97) % سيتوكروم c
					(126.5, 17.11) % ليزوزيم
					(131.5, 17.78) % GFP
					(134.0, 18.12) % GAPDH
					(137.5, 18.59) % هيموغلوبين
					(146.0, 19.74) % نيوكليوسوم
					(147.0, 19.88) % سنايسوسوم U1
					(150.5, 20.35) % ATP-سينثاز F1
					(164.0, 22.18) % ريبوسوم 70S
					(164.5, 22.24) % بروتيازوم 26S
					(187.5, 25.36) % NPC
				};
				\addlegendentry{التجمعات البروتينية}
				
				% تعليقات على الكائنات البيولوجية الرئيسية
				\node[green!60!black, font=\tiny, anchor=south west] at (axis cs:164.0,22.18) {ريبوسوم};
				\node[green!60!black, font=\tiny, anchor=south west] at (axis cs:164.5,22.24) {بروتيازوم};
				\node[green!60!black, font=\tiny, anchor=south west] at (axis cs:187.5,25.36) {NPC};
				\node[green!60!black, font=\tiny, anchor=south west] at (axis cs:122.5,16.56) {يوبيكويتين};
				\node[green!60!black, font=\tiny, anchor=south east] at (axis cs:137.5,18.59) {هيموغلوبين};
				
				% خطوط عمودية منقطة عند نصف الصحيح
				\draw[gray, thin, dotted] (axis cs:122.5,-2) -- (axis cs:122.5,30);
				\draw[gray, thin, dotted] (axis cs:125.5,-2) -- (axis cs:125.5,30);
				\draw[gray, thin, dotted] (axis cs:126.5,-2) -- (axis cs:126.5,30);
				\draw[gray, thin, dotted] (axis cs:131.5,-2) -- (axis cs:131.5,30);
				\draw[gray, thin, dotted] (axis cs:134,-2) -- (axis cs:134,30);
				\draw[gray, thin, dotted] (axis cs:137.5,-2) -- (axis cs:137.5,30);
				\draw[gray, thin, dotted] (axis cs:146,-2) -- (axis cs:146,30);
				\draw[gray, thin, dotted] (axis cs:147,-2) -- (axis cs:147,30);
				\draw[gray, thin, dotted] (axis cs:150.5,-2) -- (axis cs:150.5,30);
				\draw[gray, thin, dotted] (axis cs:164,-2) -- (axis cs:164,30);
				\draw[gray, thin, dotted] (axis cs:164.5,-2) -- (axis cs:164.5,30);
				\draw[gray, thin, dotted] (axis cs:187.5,-2) -- (axis cs:187.5,30);
			\end{axis}
		\end{tikzpicture}
		\caption{\textbf{سلم الكتلة العالمي يمتد من الجسيمات الأولية إلى أكبر الآلات الخلوية.} 
			تمثل كل نقطة اللوغاريتم الطبيعي لكتلة الجسم بالنسبة للإلكترون. 
			الخط المتقطع هو التنبؤ النظري لـ YUCT: الكتل تنمو كقوى للكم \(q = (3/2)^{1/3} \approx 1.1447\). 
			الدوائر الحمراء: فرميونات أولية؛ المربعات الزرقاء: أنوية ذرية خفيفة (لإظهار الاستمرارية)؛ المثلثات الخضراء: التجمعات البروتينية والريبونوكليوبروتينية التي نوقشت في النص. 
			جميع النقاط قريبة جداً من الخط، مما يدل على أن نفس قانون التدرج المتقطع يحكم المادة دون الذرية وفوق الجزيئية. 
			تشكل الأجسام الخضراء عنقوداً متميزاً في المدى \(N_f \approx 120\)–\(190\)، وهو الموافق لنافذة الكتلة حيث يصبح الطي الذاتي والوظيفة الخلوية المنسقة ممكنين. 
			يوفر هذا الشكل الأساس الهيكلي لساعة BCLM-EC: إذا كانت كل هذه الآلات عقداً على سلم تنسيق واحد، فإن حالتها الوظيفية – وبالتالي العلامات اللاجينية التي تبلغ عنها – يجب أن تتبع نفس قانون الخطأ العالمي.}
		\label{fig:mass_ladder_bclm}
	\end{figure}
	
	\section{تفسير الأوزان}
	\label{sec:weights}
	الوزن الموجب في الجدول~\ref{tab:cpgs} يعني أن الموقع يصبح مفرط المثيلة مع انخفاض \(\REL\) (أي مع تقدم العمر)؛ والوزن السالب يشير إلى نقص المثيلة. يعكس الحجم المطلق حساسية الموقع لحالة التنسيق. ومن الجدير بالذكر أن المواقع ذات الوزن الأعلى تقع بالقرب من الجينات المشاركة في:
	\begin{itemize}
		\item \textbf{إصلاح DNA واستقرار الجينوم} (مثل BRCA1, FANCI, RAD51C عبر NOTCH1 و TCF25)؛
		\item \textbf{وظيفة الميتوكوندريا والإجهاد التأكسدي} (مثل TXNL1, CYP2E1)؛
		\item \textbf{حفظ البروتين} (مثل HSPA13, PDCD4)؛
		\item \textbf{التصاق الخلية وإشارات المصفوفة خارج الخلية} (مثل ITGB1, COL18A1).
	\end{itemize}
	هذا الإثراء متسق مع الطبقات الأربع لبروتوكول الحياة الطويلة (الملحق BCLM)، مما يشير إلى أن انجراف المثيلة الذي تلتقطه الساعة هو مراسل لنفس الانخفاض التنسيقي الأساسي الذي يهدف البروتوكول إلى مقاومته.
	
	\section{التنبؤات وقابلية التكذيب}
	\label{sec:predictions}
	
	\begin{enumerate}
		\item \textbf{تباطؤ الشيخوخة تحت علاج الحياة الطويلة.}
		إذا رفع بروتوكول الحياة الطويلة \(\REL\) الخلوي بعامل \(f = \REL_{\mathrm{LL}}/\REL_{\mathrm{WT}}\)، فإن احتمال الخطأ ينخفض بمقدار \(f^{-\beta}\). وبالتالي، يجب أن يتباطأ انجراف المثيلة عند كل CpG في الساعة بنفس العامل. لعلاج يستمر \(t\) سنوات، ستقرأ الساعة عمراً أقل بمقدار
		\[
		\Delta\mathrm{Age} = \frac{1}{\beta\gamma} \ln f.
		\]
		باستخدام \(\beta\gamma = 0.20\) (الملحق P) وقيمة تقديرية \(f \approx 1.33\) (بناءً على تقليل الخطأ المشترك في بروتوكول الحياة الطويلة)، يتوقع \(\Delta\mathrm{Age} \approx 1.4\) سنة بعد شهر واحد من العلاج. سوف تتراكم العلاجات الأطول فترات تعويض أكبر. يمكن اختبار هذا التنبؤ بزراعة خلايا عضلة القلب المهندسة بـ LL جنباً إلى جنب مع ضوابط النمط البري ومقارنة أعمار BCLM-EC الخاصة بها على فترات منتظمة.
		
		\item \textbf{اعتماد سرعة الساعة على درجة الحرارة.}
		لأن \(\REL(T) \propto T^{-\gamma}\)، فإن الخلايا المستنبتة عند درجة حرارة أقل تمتلك \(\REL\) أعلى وبالتالي انجراف مثيلة أبطأ. على سبيل المثال، خفض درجة حرارة الاستنبات من \(37^\circ\text{C}\) إلى \(30^\circ\text{C}\) يجب أن يؤدي إلى تباطؤ الشيخوخة بحوالي 15\% على مدى عدة ممرات. هذا نتيجة مباشرة لتدرج كفاءة التنسيق مع درجة الحرارة (الملحقان BCLM و P) ويمكن اختباره في أي نوع خلايا تمت معايرة BCLM-EC له.
		
		\item \textbf{تباين الخلية الواحدة.}
		على مستوى الخلية الواحدة، في ظل افتراض أن أخطاء المثيلة مستقلة بين الخلايا، يجب أن يكون توزيع قيم المثيلة عند CpG معين في الساعة لوغاريتمياً طبيعياً مع تباين يزداد بمقدار \(\REL^{-2\beta}\). يمكن لتسلسل البايسلفايت أحادي الخلية لمجموعة شيخوخة أن يوفر بالتالي تقديراً مستقلاً لـ \(\REL\)، والذي يجب أن يتوافق مع القيمة المشتقة من مثيلة العينة المجمعة.
	\end{enumerate}
	
	\section{القيود والأسئلة المفتوحة}
	\label{sec:limitations}
	\begin{itemize}
		\item تعتمد درجة الحساسية \(S_{\mathrm{BCLM}}\) حالياً على ميل العمر \(\beta_i\) كبديل للمشتق الحقيقي \(\partial m/\partial \REL\). قياس أكثر مباشرة لـ \(\REL\) في الخلايا المفردة (على سبيل المثال، عبر التدفق الأيضي ونشاط الإصلاح) سيوفر أساساً أقوى.
		\item تم تدريب الساعة على الدم الكامل؛ يجب معايرة الساعات الخاصة بالأنسجة بشكل منفصل، لأن \(\alpha_{\mathrm{epi}}\) و \(\gamma\) قد يختلفان.
		\item نموذج الاضمحلال الأسي (\ref{eq:Kdecay}) هو تقريب من الدرجة الأولى. قد تنتج نماذج أكثر تفصيلاً لتدهور شبكة التنسيق أشكالاً غير أسية يمكن تمييزها ببيانات طولية عالية الدقة.
		\item الأوزان في الجدول~\ref{tab:cpgs} خاصة إلى حد ما بمنصة Illumina 450K. يجب تكييف الإصدارات المستقبلية مع مقايسات المثيلة القائمة على التسلسل.
	\end{itemize}
	
	\section{الاستنتاج}
	\label{sec:conclusion}
	توفر ساعة تنسيق BCLM بديلاً ميكانيكياً نظرياً للساعات اللاجينية الإحصائية البحتة. إنها متجذرة في قانون الخطأ العالمي لـ YUCT، وتختار CpGs بحساسيتها لحالة التنسيق الخلوي، وتحقق دقة تنبؤ بالعمر تنافسية مع عدد صغير من المواقع، وتقدم تنبؤات ملموسة قابلة للتكذيب حول كيفية استجابة مسارات المثيلة للتدخلات التي تعدل \(K_{\mathrm{eff}}\). إلى جانب بروتوكول الحياة الطويلة، تشكل BCLM-EC دورة تجريبية مغلقة: تتنبأ النظرية بالنتيجة اللاجينية، وتقوم النتيجة إما بالتحقق من صحة النظرية أو تكذيبها. هذا يحول دراسة الشيخوخة اللاجينية من علم وصفي إلى علم تفسيري.
	
	\vspace{0.5cm}
	\noindent\textbf{البيانات والشفرة:} جميع بيانات المثيلة متاحة من GEO (GSE40279). معاملات الساعة المدربة (الجدول~\ref{tab:cpgs})، وشيفرة R للمعايرة، وقائمة الـ 100 CpG التي تم النظر فيها أصلاً متاحة في \url{https://github.com/Alexey-Yakushev-YUCT/YPSDC}.
	
	\begin{thebibliography}{99}
		\bibitem{Horvath2013} S. Horvath, ``DNA methylation age of human tissues and cell types,'' \emph{Genome Biology}, 14, R115 (2013).
		\bibitem{Hannum2013} G. Hannum \emph{et al.}, ``Genome-wide methylation profiles reveal quantitative views of human ageing rates,'' \emph{Molecular Cell}, 49, 359--367 (2013).
		\bibitem{Levine2018} M. E. Levine \emph{et al.}, ``An epigenetic biomarker of ageing for lifespan and healthspan,'' \emph{Aging}, 10, 573--591 (2018).
		\bibitem{BCLM} A.\,V.\,Yakushev, ``Appendix BCLM: Biological Coordination Lagrangian,'' in \emph{YUCT}, Zenodo (2026).
		\bibitem{AppendixP} A.\,V.\,Yakushev, ``Appendix P: Temperature Dependence of Gravity,'' in \emph{YUCT}, Zenodo (2026).
		\bibitem{YPSDC_repo} A.\,V.\,Yakushev, YPSDC Repository, \url{https://github.com/Alexey-Yakushev-YUCT/YPSDC}.
	\end{thebibliography}
	
\end{document}